% sections/graphs.tex
\chapter{Алгоритм Дейкстры}

\section{Теоретическая часть}

Алгоритм Дейкстры~\cite{dijkstra1959} решает задачу нахождения кратчайших
путей из одной вершины-источника до всех остальных вершин во взвешенном
ориентированном графе с неотрицательными весами рёбер.

Пусть дан граф $G = (V, E, w)$, где $V$~--- множество вершин, $E$~--- множество
рёбер, $w \colon E \to \mathbb{R}_{\geq 0}$~--- весовая функция.
Требуется найти для каждой вершины $v \in V$ расстояние:

\begin{equation}
  \label{eq:dijkstra-dist}
  d(s, v) = \min_{\pi: s \leadsto v} \sum_{(u,u') \in \pi} w(u, u'),
\end{equation}

где минимум берётся по всем путям $\pi$ из источника $s$ в вершину $v$.

\subsection{Оценка сложности}

При использовании двоичной кучи (\textit{binary heap}) в качестве
приоритетной очереди сложность алгоритма составляет:

\begin{equation}
  \label{eq:dijkstra-complexity}
  T(|V|, |E|) = O\!\left((|E| + |V|)\log |V|\right).
\end{equation}

\subsection{Пример работы алгоритма}

На рисунке~\ref{fig:dijkstra-graph} изображён граф, для которого
алгоритм находит кратчайшие расстояния от вершины~\texttt{0}.

\begin{figure}[H]
  \centering
  \includegraphics[width=0.72\textwidth]{images/dijkstra_graph.png}
  \caption{Взвешенный граф и кратчайшие расстояния от вершины~\texttt{0}}
  \label{fig:dijkstra-graph}
\end{figure}

\section{Реализация на C++}

В листинге ниже приведена реализация алгоритма Дейкстры на языке~C++
с использованием приоритетной очереди из стандартной библиотеки~\cite{cormen2022}.
Граф хранится в виде списка смежности: \mintinline{cpp}{graph[u]} содержит
пары \mintinline{cpp}{(v, w)}, где \verb|v|~--- смежная вершина,
\verb|w|~--- вес ребра.

Ключевая переменная \mintinline{cpp}{dist} хранит текущие минимальные
расстояния; изначально все значения равны \mintinline{cpp}{INF}
(\verb|numeric_limits<int>::max()|), кроме расстояния до источника,
которое равно~\texttt{0}.

\inputminted[
  label=\fbox{dijkstra.cpp}
]{cpp}{code/dijkstra.cpp}
